% ==============================================================================
% 02_hien_trang.tex — Phần 2. Khảo sát hiện trạng (AS-IS)
% Phụ trách: Thành viên 3
% Ngân sách trang: 3 trang
% ==============================================================================

\pttkPart{Khảo sát hiện trạng (AS-IS)}

\todosection{Thành viên 3}{Số liệu khảo sát phải nhất quán với Phần 1 (OKR) và Phụ lục A, B. Ngân sách mục tiêu: 3 trang.}

\subsection{Đặc điểm các nhóm người dùng khảo sát}
% Bảng 3 phân khúc mục tiêu: Solo Dev, Nhóm nhỏ (Sinh viên/Startup), Nhóm nhiều module.

\subsection{Phương pháp và kế hoạch khảo sát}
% Bảng kế hoạch & kết quả khảo sát (Phỏng vấn sâu, Bảng hỏi trực tuyến, Quan sát thực tế, Phân tích repo/log CI).

\subsection{Quy trình quản lý công việc và xử lý lỗi hiện tại}
\todofig{Hình 1. Quy trình quản lý công việc và xử lý lỗi CI hiện tại (AS-IS) — diagrams/out/hinh\_01\_asis.pdf}

\subsection{Hệ thống và dữ liệu hiện có}
% Liệt kê công cụ rời rạc: Trello/Jira/GitHub Issues, GitHub/GitLab, CI logs, IDE (Cursor/VS Code).

\subsection{Điểm đau (Pain points)}
% Bảng danh sách 6 điểm đau (P1..P6) kèm ảnh hưởng và bằng chứng định lượng.
